% BHU PYQ -- Transcribed & Corrected
% Paper: CHEMJ11: Basic Concepts of Chemistry-I
% Exam: B.Sc. (Hons) Semester I (NEP) Examination 2024-25
% Ref Code: CECONF/N/2024-25/28
% Category: Major
% Department: Chemistry

\documentclass[11pt,a4paper]{article}
\usepackage[a4paper,top=2.5cm,bottom=2.5cm,left=2.8cm,right=2.8cm,headheight=14pt]{geometry}
\usepackage{amsmath,amssymb}
\usepackage[T1]{fontenc}
\usepackage[utf8]{inputenc}
\usepackage{lmodern,microtype,fancyhdr,enumitem,hyperref,lastpage,parskip}

\hypersetup{
    colorlinks=true,
    urlcolor=black,
    linkcolor=black,
    pdftitle={CHEMJ11: Basic Concepts of Chemistry-I -- BHU Sem I 2024-25 (NEP)}
}

\pagestyle{fancy}
\fancyhf{}
\renewcommand{\headrulewidth}{0.4pt}
\renewcommand{\footrulewidth}{0.4pt}
\fancyhead[L]{\small \textit{Banaras Hindu University}}
\fancyhead[R]{\small \textit{CHEMJ11: Basic Concepts of Chemistry-I (2024-25)}}
\fancyfoot[C]{\small Page \thepage\ of \pageref{LastPage}}

\newcommand{\pts}[1]{\hfill \textit{[#1]}}

\begin{document}

\begin{center}
  {\large\bfseries BANARAS HINDU UNIVERSITY}\\[2pt]
  {\normalsize B.Sc. (Hons) Semester I Examination, 2024-25 (NEP)}\\[6pt]
  \rule{0.8\linewidth}{0.5pt}\\[6pt]
  {\large\bfseries CHEMISTRY}\\[4pt]
  {\large\bfseries Paper Code: CHEMJ11 : Basic Concepts of Chemistry-I}\\[2pt]
  {\small Ref. No.: CECONF/N/2024-25/28}\\[4pt]
  {\normalsize Time: 2$\frac{1}{2}$ Hours\hfill Maximum Marks: 70}
\end{center}

\rule{\linewidth}{0.4pt}

\smallskip

\noindent \textit{Instructions: Attempt five questions including Question 1 which is compulsory. Attempt at least one question from Questions 2 and 3. The figures in the right-hand margin indicate marks.}

\bigskip

\noindent \textbf{Question 1.}
\begin{enumerate}[label=(\alph*)]
  \item 
  \begin{enumerate}[label=(\i*)]
    \item What is the physical significance of + and - signs in the orbitals? \pts{1}
    \item What are the conditions for the acceptability of wave functions obtained by solving the Schrödinger wave equation? \pts{2}
  \end{enumerate}
  \item What are the different consequences of the inert pair effect? \pts{2}
  \item Obtain the expressions for van der Waals constants $a$ and $b$ in terms of $T_c$, $P_c$ and $R$. \pts{3}
  \item State the first law of thermodynamics and write its limitations. \pts{3}
  \item Prove that the work done in an isothermal case is higher than in an adiabatic reversible expansion of an ideal gas. \pts{3}
\end{enumerate}

\bigskip

\noindent \textbf{Question 2.}
\begin{enumerate}[label=(\alph*)]
  \item Calculate the effective nuclear charge on the outermost electron of $^{29}\text{Cu}$. \pts{2}
  \item What are eigenfunctions and eigenvalues? Give two examples for each. \pts{2}
  \item What are the trends of the solubilities of alkali-metal fluorides and iodides in water? Discuss the reasons for the solubility trends of these salts. \pts{4}
  \item What are the limitations of Born-Landé equation? \pts{2}
  \item 
  \begin{enumerate}[label=(\i*)]
    \item Calculate the standard reduction potential for $\text{Fe}^{3+}/\text{Fe}$. Given $\text{Fe}^{3+}/\text{Fe}^{2+} = +0.77\text{ V}$ and $\text{Fe}^{2+}/\text{Fe} = -0.47\text{ V}$. \pts{2}
    \item Will $\text{Fe}^{2+}$ disproportionate to $\text{Fe}^{3+}$ and $\text{Fe}$ in an acid solution? Support your answer with reasons. \pts{4}
  \end{enumerate}
\end{enumerate}

\bigskip

\noindent \textbf{Question 3.}
\begin{enumerate}[label=(\alph*)]
  \item Obtain the Schrödinger wave equation assuming the electron behaves like a standing wave in three dimensions. \pts{4}
  \item Qualitatively discuss the trends of 1st ionization energies for the elements in the 2nd and 3rd rows of the periodic table. \pts{3}
  \item Discuss the dependence of lattice energy in terms of ionic size and explain the consequences of ionic size on lattice energy and hydration energy. \pts{3}
  \item Calculate the electronegativity of carbon by Pauling's method if $\Delta$, the extra bond energy for the C-H bond is $24.3\text{ kJ\,mol}^{-1}$. Explain, what is meant by extra bond energy. \pts{4}
\end{enumerate}

\bigskip

\noindent \textbf{Question 4.}
\begin{enumerate}[label=(\alph*)]
  \item Write the basic postulates of kinetic molecular theory of gases giving proper evidences for all of them. \pts{4}
  \item Explain the terms $T_c$, $P_c$ and $V_c$ with suitable examples. How will you determine the value of $V_c$ experimentally? \pts{5}
  \item State the law of corresponding states and derive the expression for van der Waals reduced equation of state. \pts{5}
\end{enumerate}

\bigskip

\noindent \textbf{Question 5.}
\begin{enumerate}[label=(\alph*)]
  \item The number of drops of liquid and water was observed to be 66 and 30 respectively by the number drop method using a stalagmometer. Calculate the surface tension of the liquid. The surface tension of water is $72.75\text{ dyne/cm}$ and masses of liquid and water are $10.750\text{ g}$ and $11.250\text{ g}$, respectively for the same volume. \pts{4}
  \item What is the law of rational indices? What is the need to convert Weiss indices into Miller indices? Write the steps involved. \pts{5}
  \item Draw and explain the crystal structures of graphite and diamond. \pts{5}
\end{enumerate}

\bigskip

\noindent \textbf{Question 6.}
\begin{enumerate}[label=(\alph*)]
  \item Using $\mu_{JT} = -\frac{1}{C_p} \left[ \left(\frac{dE}{dV}\right)_T \cdot \left(\frac{dV}{dP}\right)_T \right] - \frac{1}{C_p} \left[ \frac{d(PV)}{dP} \right]_T$, show under what conditions the Joule-Thomson coefficient will be positive for real gas. \pts{4}
  \item Why cooling of real gas is observed in Joule-Thomson effect? Show that the Joule-Thomson effect is an isenthalpic process. \pts{5}
  \item What is the physical significance of entropy? ``The entropy of the universe is always increasing.'' Justify the statement. \pts{5}
\end{enumerate}

\bigskip

\noindent \textbf{Question 7.}
\begin{enumerate}[label=(\alph*)]
  \item Prove the Maxwell relation $\left(\frac{dP}{dT}\right)_V = \left(\frac{dS}{dV}\right)_T$. \pts{4}
  \item The total heat absorbed by the system cannot be converted into useful work. Explain the statement based on the Carnot cycle. \pts{5}
  \item Show that $-(\Delta G)_T = W_{\text{Network}}$. Under what condition the value of Gibbs free energy and work function would be equal? \pts{5}
\end{enumerate}

\bigskip
\begin{center}
  \textbf{***}
\end{center}

\end{document}
